\documentclass[a4paper,11pt]{article}

%% Basics
\usepackage[margin=2cm]{geometry}
\usepackage{graphicx}
\usepackage{hyperref}
\usepackage{xltxtra}
\usepackage{xgreek}
\setmainfont[Mapping=tex-text]{Linux Libertine O}

%% Extras
\usepackage{enumerate}
\usepackage{fancyhdr}
\usepackage{fontawesome}
\usepackage{titling}

\pagestyle{fancy}
\lfoot{\faInstagram/\faFacebook \texttt{Footvolley Athens}}
\rfoot{\faEnvelopeO \texttt{info@futevolei.gr}}
\renewcommand{\chaptername}{}
\renewcommand{\headrulewidth}{0pt}
\renewcommand{\footrulewidth}{0.4pt}

\pretitle{%
  \begin{center}
  \LARGE
  \includegraphics[
    width=0.3\textwidth,
    keepaspectratio
  ]{../images/Footvolley-logo-cropped-reduced.pdf}\\[\bigskipamount]
}
\posttitle{\end{center}}
%\setlength{\droptitle}{-1cm}

\title{Open Footvolley Tournament \\ Κανονισμοί Διοργάνωσης}
\date{}

\begin{document}

\maketitle
\thispagestyle{fancy}

\vspace{-2cm}

\begin{enumerate}

\item Το τεχνικό meeting ξεκινάει στις 10:00 και όλοι/ες οι παίκτες/ριες πρέπει
  να είναι παρόντες/ούσες.

\item Στα παιχνίδια χρησιμοποιούνται \textbf{αποκλειστικά} μπάλες του Footvolley
  Athens.

\item To σύστημα των αγώνων θα είναι όμιλοι στην α' φάση και knock-out από τους
  προημιτελικούς/ημιτελικούς. Aν ο αριθμός των ομάδων είναι μεγάλος, θα γίνει
  \href{https://www.printyourbrackets.com/9teamdoubleelimination.html}{ταμπλό με
    διπλό αποκλεισμό}.

\item Όλα τα παιχνίδια αποτελούνται από ένα σετ στους 18 πόντους.

\item Για τη κατάταξη των ομάδων στον όμιλο παίζουν ρόλο τα εξής κριτήρια με τη
  σειρά που αναγράφονται:

  \begin{enumerate}[i)]
  \item Αριθμός νικών
  \item Αριθμός νικών στην/ις αναμέτρηση/σεις μεταξύ των ισοβαθμούντων ομάδων
  \item Συνολικοί πόντοι υπέρ μείον συνολικοί πόντοι κατά (max.) στις
    αναμετρήσεις μεταξύ των ισοβαθμούντων ομάδων
  \item Συνολικοί πόντοι υπέρ (max.) στις αναμετρήσεις μεταξύ των ισοβαθμούντων
    ομάδων
  \item Κλήρωση %ή σετ στα 12 (max. στα 15)
  \end{enumerate}

\item Η ομάδα που έχασε το τελευταίο παιχνίδι διαιτητεύει το επόμενο παιχνίδι.
  Μετά το τέλος του παιχνιδιού, οι διαιτητές ενημερώνουν τη γραμματεία για το
  ακριβές σκορ.

\item Στην αρχή του παιχνιδιού αποφασίζεται ποια ομάδα θα έχει το πρώτο service
  με ρίψη κέρματος. Το ίδιο συμβαίνει και στην περίπτωση που χρειαστεί
  tie-break.

\item Οι παίκτες μπορούν να χρησιμοποιήσουν οποιοδήποτε μέρος του σώματος εκτός
  από τα χέρια και τα μπράτσα.

\item Η κάθε ομάδα δικαιούται να περάσει την μπάλα χρησιμοποιώντας 1 εώς 3
  επαφές.

\item Απαγορεύεται κάποιος παίκτης να κάνει δύο \textit{συνεχόμενες} επαφές.

\item Σε \textbf{καμία} περίπτωση δεν επιτρέπεται να ακουμπήσει κάποιος από τους
  παίκτες το φιλέ. Αυτό αποτελεί παράβαση και ο διαιτητής κατοχυρώνει τον πόντο
  στην αντίπαλη ομάδα.

\item Ο παίκτης που επιτίθεται ή αμύνεται μπορεί να περάσει στην αντίπαλη
  περιοχή χωρίς να ακουμπήσει το φιλέ ή να εμποδίσει την αντίπαλη ομάδα.

\item Η μπάλα πρέπει να περνάει στην απέναντι περιοχή πάνω από το φιλέ και
  ανάμεσα των δύο κεραιών ή των νοητών προεκτάσεων αυτών. Αν η μπάλα περάσει
  εκτός αυτών ή ακουμπήσει κάποια εξ αυτών ή την περιοχή του φιλέ εκτός αυτών,
  καταλογίζεται παράβαση και η αντίπαλη ομάδα κερδίζει τον πόντο.

\item Η ζώνη του service ορίζεται από την πίσω γραμμή της περιοχής εώς και 1m
  πίσω από αυτή. Αν το βουνό του service καλύπτει έστω και λίγο τη γραμμή,
  \textit{ολόκληρο} το βουνό αποτελεί κομμάτι της περιοχής της ομάδας.

\item Κατά τη διάρκεια του service, ο συμπαίκτης δεν μπορεί να εμποδίζει τους
  παίκτες της αντίπαλης ομάδας να δουν την μπάλα (π.χ. κουνώντας τα χέρια του,
  πηδώντας ή κρύβοντας την πορεία της μπάλας).

\item Αν μπει μια εξωτερική μπάλα στο γήπεδο κατά τη διάρκεια του παιχνιδιού, ο
  διαιτητής οφείλει να διακόψει το παιχνίδι και να επαναληφθεί το service.

\item Το time-out διαρκεί 1 λεπτό και η κάθε ομάδα έχει ένα time-out ανά σετ.

\item Αν μία ομάδα δεν είναι έτοιμη να αγωνιστεί εντός 5 λεπτών από την ώρα που
  θα κληθεί στο γήπεδο, χάνει το σετ 15x0, 18x0 ή 21x0 (αναλόγως από το σύστημα
  των αγώνων, βλ. προηγούμενο κανόνα).

\item Δεν μπορεί να αντικατασταθεί παίκτης από κάποια ομάδα από τη στιγμή που θα
  ξεκινήσουν οι αγώνες.

\item Η οργανωτική επιτροπή διατηρεί το δικαίωμα να απορρίψει ή να αποκλείσει
  οποιοδήποτε παίκτη ή ομάδα που δημιουργεί προβλήματα στην ομαλή διεξαγωγή των
  αγώνων. Τυχούσες ελλείψεις ή αμφισβητούμενες περιπτώσεις του παρόντος
  κανονισμού θα αναλύονται και θα λύνονται από την οργανωτική επιτροπή.

\end{enumerate}

%\renewcommand{\section}[2]{}% Hide section name "Αναφορές"
\renewcommand{\refname}{}

\begin{thebibliography}{1}

\bibitem[EFVL]{EFVL}
  European Footvolley League:
  \textit{Official Footvolley Rules}
  \\\texttt{https://footvolleyeurope.com/official-footvolley-rules/}

\bibitem[CGF]{CGF}
  Circuito Gaúcho de Futevôlei:
  \textit{Regulamento}
  \\\texttt{https://www.circuitogauchodefutevolei.com/regulamento}
\end{thebibliography}

\end{document}
